\section*{Chapitre \ref{ch:b2:catalytic-cycles} --- Catalyse organométallique~: étapes élémentaires et cycles}

\begin{solution}{exo:b2:catalytic-cycles:1}
Avant~: Ir(I), $d^8$, $8 + 4 \times 2 = 16$ électrons. Après~: Ir(III), $d^6$, $6 + 6 \times 2 = 18$
électrons, octaédrique.
\end{solution}

\begin{solution}{exo:b2:catalytic-cycles:2}
La première est une \omterm{def:b2:catalytic-cycles:oxidative-addition}{élimination réductrice} (Pt(IV) en Pt(II), éthane formé à partir de deux méthyles).
La seconde est un \omterm{def:b2:catalytic-cycles:ligand-exchange}{échange de ligand} (une phosphine remplacée par l'éthène) suivi d'une
\omterm{def:b2:catalytic-cycles:insertion}{insertion migratoire} de l'éthène dans la liaison Pd--Ph.
\end{solution}

\begin{solution}{exo:b2:catalytic-cycles:3}
TON $= 15.0/0.020 = 750$~; TOF $= 750/3.0 = \qty{250}{h^{-1}}$.
\end{solution}

\begin{solution}{exo:b2:catalytic-cycles:4}
\begin{align*}
  &\mathrm{PdL_2 + ArBr \to ArPdBrL_2}\\
  &\mathrm{ArPdBrL_2 + Ar'B(OH)_3^- \to ArPdAr'L_2 + Br^- + B(OH)_3}\\
  &\mathrm{ArPdAr'L_2 \to Ar{-}Ar' + PdL_2}
\end{align*}
Somme~: $\mathrm{ArBr + Ar'B(OH)_3^- \to Ar{-}Ar' + Br^- + B(OH)_3}$~; toutes les espèces du palladium
se simplifient.
\end{solution}

\begin{solution}{exo:b2:catalytic-cycles:5}
L'éthyle et l'isopropyle (ils portent des hydrogènes $\beta$). Le méthyle n'a pas de carbone $\beta$~; le carbone $\beta$
du néopentyle ne porte aucun hydrogène~; le carbone $\beta$ du benzyle est un carbone \omterm{def:b2:huckel:aromatic}{aromatique} sans hydrogène, et aucun hydrogène du cycle ne peut atteindre
le métal.
\end{solution}

\begin{solution}{exo:b2:catalytic-cycles:6}
(\textit{E})-3-phénylpropénoate de méthyle, \ce{Ph-CH=CH-COOCH3}, le phényle étant ajouté sur le
carbone terminal~; produit trans.
\end{solution}

\begin{solution}{exo:b2:catalytic-cycles:7}
Le butanal \ce{CH3CH2CH2CHO} (linéaire~: le métal se fixe sur le carbone terminal, l'hydrure sur
le carbone interne) et le 2-méthylpropanal \ce{(CH3)2CHCHO} (ramifié~: le métal sur le carbone
interne).
\end{solution}

\begin{solution}{exo:b2:catalytic-cycles:8}
\ce{[RhCl(PPh3)3]} a 16 électrons~: additionner \ce{H2} (+2) donnerait 18 électrons avec six
positions de coordination, trop encombré avec trois phosphines volumineuses. \ce{[RhCl(PPh3)2]} en a 14
et possède un site libre~: l'\omterm{def:b2:catalytic-cycles:oxidative-addition}{addition oxydante} donne facilement un complexe à 16 électrons.
\end{solution}

\begin{solution}{exo:b2:catalytic-cycles:9}
L'acide boronique est un mauvais nucléophile~; la base ajoute un hydroxyde (ou un alcoolate) au bore,
ce qui donne un borate \ce{ArB(OH)3-}, dont le groupe aryle est plus riche en électrons et passe sur le
palladium.
\end{solution}

\begin{solution}{exo:b2:catalytic-cycles:10}
Un complexe à 20 électrons n'est pas raisonnable. Le complexe doit d'abord perdre un ligand (14 e), puis
additionner \ce{H2} (16 e), fixer l'alcène (18 e), l'insérer dans une liaison M--H (16 e) et éliminer
l'alcane (14 e), avant de reprendre le ligand ou de recommencer.
\end{solution}

\begin{solution}{exo:b2:catalytic-cycles:11}
Vitesse initiale $n_{\max}/\tau = \qty{0.25}{mmol/min}$~; TOF initiale $= 0.25/0.010 =
\qty{25}{min^{-1}} = \qty{1.5e3}{h^{-1}}$. TON final $= 10.0/0.010 = 1.0 \times 10^3$.
\end{solution}

\begin{solution}{exo:b2:catalytic-cycles:12}
Le 4-bromoanisole avec l'acide phénylboronique (ou le bromobenzène avec l'acide 4-méthoxy\-phényl\-boronique),
un catalyseur phosphine--palladium(0) et une base. Cycle~: \omterm{def:b2:catalytic-cycles:oxidative-addition}{addition oxydante} du
bromure d'aryle, \omterm{def:b2:catalytic-cycles:transmetalation}{transmétallation} à partir du borate, \omterm{def:b2:catalytic-cycles:oxidative-addition}{élimination réductrice} du
4-méthoxybiphényle.
\end{solution}

\begin{solution}{pb:b2:catalytic-cycles:1}
\textbf{1.} Pd(II), $d^8$.
\textbf{2.} Pd(0), $d^{10}$, $10 + 2 \times 2 = 14$ électrons.
\textbf{3.} Avec 14 électrons, il peut accepter deux doublets de plus~: il possède des \omterm{def:b2:catalytic-cycles:unsaturated}{sites vacants}.
\textbf{4.} Les complexes phosphine--palladium(0) sont oxydés par l'air, et les phosphines sont
oxydées en oxydes de phosphine.
\textbf{5.} Elle réduit le palladium(II) en palladium(0) (en s'oxydant elle-même) et se lie au
palladium(0), qu'elle maintient en solution.
\textbf{6.} Le \omterm{def:b2:catalytic-cycles:cycle}{précatalyseur}.
\textbf{7.} \ce{Pd(PPh3)2 + CH3C6H4Br -> [Pd(C6H4CH3)Br(PPh3)2]}~: Pd(II), $d^8$, $8 + 4 \times 2 =
16$ électrons.
\textbf{8.} Il le transforme en borate \ce{PhB(OH)3-}, qui transfère son groupe phényle.
\textbf{9.} \ce{[Pd(Ar)Br(PPh3)2] + PhB(OH)3- -> [Pd(Ar)(Ph)(PPh3)2] + Br- + B(OH)3}.
\textbf{10.} L'\omterm{def:b2:catalytic-cycles:oxidative-addition}{élimination réductrice} unit deux ligands en cis~; des aryles en trans doivent d'abord s'isomériser.
\textbf{11.} \ce{[Pd(Ar)(Ph)(PPh3)2] -> Ar-Ph + Pd(PPh3)2}~: palladium(0), 14 électrons, prêt
pour un nouveau tour.
\textbf{12.} \ce{CH3C6H4Br + PhB(OH)3- -> CH3C6H4-Ph + Br- + B(OH)3}.
\textbf{13.} Le 4-bromotoluène (\qty{10.0}{mmol}, contre \qty{12.0}{mmol} d'acide boronique).
\textbf{14.} $1.51/168.24 = \qty{8.98e-3}{mol}$, \qty{8.98}{mmol}.
\textbf{15.} $8.98/10.0 = 89.8~\%$.
\textbf{16.} $0.050/10.0 = 0.50$~\%~molaire.
\textbf{17.} $8.98/0.050 = 180$.
\textbf{18.} $180/4.0 = \qty{45}{h^{-1}}$.
\textbf{19.} Deux groupes phényle transférés sur le même palladium (homocouplage de
l'acide boronique, favorisé par des traces de dioxygène), puis éliminés ensemble.
\textbf{20.} Les groupes aryle n'ont pas d'hydrogène $\beta$ capable d'atteindre le métal.
\textbf{21.} L'\omterm{def:b2:catalytic-cycles:oxidative-addition}{addition oxydante}~: la liaison \ce{C-Cl} est plus forte que \ce{C-Br}.
\textbf{22.} Il finit sous forme de particules de palladium(0) ou de sels dans les résidus~; il est cher et
toxique, et on le récupère.
\textbf{23.} Une partie de l'acide boronique est perdue dans des réactions parasites (homocouplage, perte du bore
du cycle)~; l'excès maintient le bromure limitant.
\textbf{24.} Le palladium a accompli $\boldsymbol{\approx \num{1.8e2}}$ cycles.
\end{solution}
